\section{买卖合同}
\begin{enumerate}
  \item 双务、有名、有偿、诺成、不要式。必考客观、主观。
        \begin{enumerate}
          \item 无权处分。
                \begin{enumerate}
                  \item 物债两分，区分原则。不影响效力只影响履行。
                  \item 解除合同 + 违约责任。
                \end{enumerate}
          \item 出卖人义务。
                \begin{enumerate}
                  \item 主给付义务。交付。
                  \item \zd{从给付义务}。单证资料、配套设施。
                        \begin{enumerate}
                          \item 违反构成违约 \yjt 违约责任。
                          \item 违反原则不可以解除合同、不影响风险转移、不构成重大误解、不得主张惩罚性赔偿。
                        \end{enumerate}
                  \item 回收义务。
                  \item 附随义务。通知、协助等。
                \end{enumerate}
          \item 买受人。\zd{多履行不构成违约}、及时通知。
          \item 风险负担（价金风险）。\zd{利益控制理论。利益风险同一，谁享受谁承担}。有约从约。
                \begin{enumerate}
                  \item 风险 = \zd{标的灭失 + 不是双方过错}。不可抗力、第三人原因。
                  \item \zd{交付主义}。交付 = 移转（主客观相结合。移转标的物所有权的意思 + 占有）。
                        \begin{enumerate}
                          \item 不动产。分离主义。房屋 + 买卖。
                                \begin{itemize}
                                  \item 交付 √（债）+ 登记 ×（物）\yjt 风险转移，有债无物。
                                  \item 交付 × + 登记 √ \yjt 风险未转移，所有权转移，有物无债。
                                  \item 交付 √ + 登记 √ \yjt 风险转移、所有权转移，有债有物。
                                \end{itemize}
                          \item 动产。合一主义。考例外。
                                \begin{itemize}
                                  \item 动产所有权保留 \yjt 所有权和风险交付分离。
                                  \item 无约定。\zd{种类物特定化}（\al{已装箱}）\hz 装运单据 \hz 加盖标记 \hz 通知买受人 \yjt 视为转移。
                                \end{itemize}
                        \end{enumerate}
                  \item （根本）违约在先。谁违约谁承担。重点考房屋买卖过程中，\zd{主体结构质量不合格}。
                  \item 代办托运（委托）、在途货物（单证买卖）。结合一起考。
                        \begin{enumerate}
                          \item 交付地点不确定 \yjt 货交第一承运人。
                          \item 在途货物 \yjt 合同成立时（因为很难确定标的物发生毁损灭失的确切时间点）。
                        \end{enumerate}
                \end{enumerate}
                % https://q.uiver.app/?q=WzAsMTEsWzEsMCwi5YyX5LqsXFxcXOeUslxcXFzljZYiXSxbNiwwLCLlub/lt55cXFxc5LmZXFxcXOS5sCJdLFswLDIsIuS4gSJdLFsyLDIsIuS4mSJdLFs1LDEsIuW5v+S4nCJdLFs2LDEsIuW5v+ilvyJdLFs3LDEsIuS6keWNlyJdLFs4LDEsIui0teW3niJdLFswLDAsIuWPl+aJmOS6uiIsWzEyMCw2MCw2MCwxXV0sWzcsMCwi5aeU5omY5Lq6IixbMTIwLDYwLDYwLDFdXSxbOCwwLCIyMDIyLjExLjE4XFxcXOaIkOeri+aXtiIsWzI3MCw2MCw2MCwxXV0sWzAsMiwi5L+d6Zmp5ZCI5ZCMIiwxXSxbMCwzLCLov5DovpPlkIjlkIwiLDFdLFsxLDQsIjEuMlciLDFdLFsxLDUsIjEuMlciLDFdLFsxLDYsIjEuMVciLDFdLFsxLDcsIjEuM1ciLDFdLFswLDEsIjIwMCDlj7DnlLXohJEiXSxbMiwwLCLkv53ljZUiLDEseyJjdXJ2ZSI6LTMsImNvbG91ciI6WzAsNjAsNjBdfSxbMCw2MCw2MCwxXV0sWzMsMCwi5o+Q5Y2VIiwxLHsiY3VydmUiOjMsImNvbG91ciI6WzAsNjAsNjBdfSxbMCw2MCw2MCwxXV0sWzEsMCwi5Luj5Yqe5omY6L+QL+WnlOaJmOWQiOWQjCIsMCx7InN0eWxlIjp7ImhlYWQiOnsibmFtZSI6Im5vbmUifX19LFsxMjAsNjAsNjAsMV1dLFszLDQsIuS6pOS7mOaXtlxcXFwyMDIyLjExLjE2IiwwLHsiY29sb3VyIjpbMjcwLDYwLDYwXX0sWzI3MCw2MCw2MCwxXV0sWzMsNSwiIiwxLHsiY29sb3VyIjpbMjcwLDYwLDYwXX1dLFszLDYsIiIsMSx7ImNvbG91ciI6WzI3MCw2MCw2MF19XSxbMyw3LCLlnKjpgJTotKfniakiLDIseyJjb2xvdXIiOlsyNzAsNjAsNjBdfSxbMjcwLDYwLDYwLDFdXV0=
\[
  \begin{tikzcd}[mytest]
    \textcolor{mygreen}{\txt{受托人}} & \txt{北京\\甲\\卖} &&&&& \txt{广州\\乙\\买} & \textcolor{mygreen}{\txt{委托人}} & \textcolor{magenta}{\txt{成立时\\2022.1.8}} \\
    &&&&& \txt{广东} & \txt{广西} & \txt{云南} & \txt{贵州} \\
    \txt{丁} && \txt{丙}
    \arrow["\txt{保险合同}"{description}, from=1-2, to=3-1]
    \arrow["\txt{运输合同}"{description}, from=1-2, to=3-3]
    \arrow["\txt{1.2W}"{description}, from=1-7, to=2-6]
    \arrow["\txt{1.2W}"{description}, from=1-7, to=2-7]
    \arrow["\txt{1.1W}"{description}, from=1-7, to=2-8]
    \arrow["\txt{1.3W}"{description}, from=1-7, to=2-9]
    \arrow["\txt{200 台电脑 1W/台}", from=1-2, to=1-7]
    \arrow["\txt{保单}"{description}, color={red}, curve={height=-30pt}, from=3-1, to=1-2]
    \arrow["\txt{提单}"{description}, color={red}, curve={height=30pt}, from=3-3, to=1-2]
    \arrow["\txt{代办托运/委托合同}"{text={mygreen}}, no head, from=1-7, to=1-2]
    \arrow["\txt{交付时\\2022.1.6}", color={magenta}, from=3-3, to=2-6]
    \arrow[color={magenta}, from=3-3, to=2-7]
    \arrow[color={magenta}, from=3-3, to=2-8]
    \arrow["\txt{在途货物}"', color={magenta}, from=3-3, to=2-9]
  \end{tikzcd}
\]

          \item 孳息归属。有约从约 \yjt 交付主义。
        \end{enumerate}
  \item 3 类特种买卖合同。
        \begin{enumerate}
          \item 试用买卖。\zd{附生效条件的买卖合同}。
                \begin{enumerate}
                  \item 生效条件。买。
                  \item 试用期的确定。
                        \begin{enumerate}
                          \item 约定 \yjt 出卖人确定。
                          \item 沉默、积极的行为视为购买。\zd{转卖}（先占有后购买，不是无权处分）。
                        \end{enumerate}
                  \item 风险。卖方。没考过。
                \end{enumerate}
          \item 动产所有权保留（\zd{本质是担保}）。只能是动产，保护的是卖家。考的最多。
                \begin{enumerate}
                  \item 价款履行问题。\zd{< 75\% \yjt 卖方取回权}；补交剩余 \yjt 买方回赎权。
                  \item 登记对抗。\zd{交付不转移所有权}。没考过。
                  \item vs 购买价金担保权。用哪个，看买受人支付能力、意愿。
                \end{enumerate}
          \item 分期付款。提前消费问题。
                \begin{enumerate}
                  \item 价款给付方式。\zd{至少分 3 次 + 80\%}。
                  \item 两线三区。动产所有权保留 + 分期付款一起考。
                  \item 特殊行权。卖方，三权择一。
                        \begin{enumerate}
                          \item \zd{加速到期权/变更权}。丧失分期付款权利，剩余价款一次性付清。
                          \item 解除权。根本违约。恢复原状 \yjt 给付型的不当得利，双方互相返还。伤敌一千自损八百。
                          \item \zd{取回权}。还能一毛不还回去。一般选择行使取回权。
                          \item 违约责任。
                        \end{enumerate}
                  \item 登记对抗。未经登记不得对抗善意第三人。
                \end{enumerate}
        \end{enumerate}
  \item 一物数卖。
        \begin{enumerate}
          \item 不动产。数卖没有顺序限制，数租有。
          \item 动产。
                \begin{enumerate}
                  \item 普通动产。受领交付 > 支付价款 > 合同成立顺序。
                  \item 特殊动产（船舶、航空器、机动车）。受领交付 > 登记 > 合同成立顺序。
                \end{enumerate}
        \end{enumerate}
\end{enumerate}
